\subsubsection*{Solution}
\paragraph*{Step-by-Step Derivation}

Assume:

\begin{itemize}[leftmargin=*]
\item Both input sidebands and all loss ports are in the vacuum state.
\item $\mu$ and $\eta$ are power transmissivities, so the corresponding field-amplitude factors are $\sqrt{\mu}$ and $\sqrt{\eta}$.
\item The photocurrent is normalized so that vacuum noise equals unity:
\end{itemize}
  \[
\left\langle |I_\theta(\nu)|^2\right\rangle_{\mathrm{vac}}=1.
\]

The two-mode Bogoliubov transformation used below is the standard quantum description of a parametric amplifier, satisfying $|\cosh r|^2-|\sinh r|^2=1$ \href{\detokenize{https://ocw.mit.edu/courses/6-453-quantum-optical-communication-fall-2016/f33f904ae6ac106c54ee07d21ab4cc63_MIT6_453F16_Lect12_Notes.pdf}}{MIT quantum-optics notes}. Optical loss is modeled as a beam splitter that introduces an independent vacuum field with amplitude coefficient $\sqrt{1-\eta}$ \href{\detokenize{https://www.nature.com/articles/s41534-017-0020-8}}{npj Quantum Information}.

\subparagraph*{1. Sideband-vector representation}

Let

\[
a_s\equiv a_{\Omega+\nu},\qquad a_i\equiv a_{\Omega-\nu},
\]

and introduce

\[
\mathbf{a}=
\begin{pmatrix}
a_s\\
a_i^\dagger
\end{pmatrix}.
\]

The photocurrent sideband is

\[
I_\theta(\nu)
=
\begin{pmatrix}
e^{-i\theta}&e^{i\theta}
\end{pmatrix}
\mathbf a
\equiv \mathbf q_\theta\mathbf a.
\]

Since

\[
[I_\theta(\nu),I_\theta^\dagger(\nu)] = [a_s,a_s^\dagger]+[a_i^\dagger,a_i] =1-1=0,
\]

operator ordering does not affect the measured spectral power. Thus,

\[
\left\langle |I_\theta(\nu)|^2\right\rangle = \left\langle I_\theta(\nu)I_\theta^\dagger(\nu)\right\rangle.
\]

For a vacuum sideband pair,

\[
K_0\equiv
\left\langle \mathbf a_0\mathbf a_0^\dagger\right\rangle
=
\begin{pmatrix}
1&0\\
0&0
\end{pmatrix}.
\]

\subparagraph*{2. OPA transformation}

Choose the pump-phase convention

\[
U(r,\phi)=
\begin{pmatrix}
\cosh r&e^{i\phi}\sinh r\\
e^{-i\phi}\sinh r&\cosh r
\end{pmatrix}.
\]

Thus,

\[
\mathbf a_{\mathrm{out}}=U(r,\phi)\mathbf a_{\mathrm{in}}.
\]

Write

\[
U_j\equiv U(r_j,\phi_j),\qquad c_j=\cosh r_j,\qquad s_j=\sinh r_j.
\]

\subparagraph*{3. Include on-chip and detection losses}

Let $\mathbf v$ be the vacuum coupled in by the on-chip loss and $\mathbf w$ the vacuum coupled in by detection loss. The complete input-output relation is

\[
\mathbf a_d = \sqrt{\eta}\, U_2\left( \sqrt{\mu}\,U_1\mathbf a_0 + \sqrt{1-\mu}\,\mathbf v \right) + \sqrt{1-\eta}\,\mathbf w.
\]

Because $\mathbf a_0$, $\mathbf v$, and $\mathbf w$ are independent vacuum fields, all cross correlations between them vanish. Therefore,

\[
\adjustbox{max width=\linewidth}{$\displaystyle \begin{aligned}
\left\langle |I_\theta(\nu)|^2\right\rangle
={}&
\eta\mu\,
\mathbf q_\theta U_2U_1K_0
(U_2U_1)^\dagger\mathbf q_\theta^\dagger\\
&+\eta(1-\mu)\,
\mathbf q_\theta U_2K_0U_2^\dagger
\mathbf q_\theta^\dagger\\
&+(1-\eta)\,
\mathbf q_\theta K_0\mathbf q_\theta^\dagger.
\end{aligned}$}
\]

The last term is simply

\[
\mathbf q_\theta K_0\mathbf q_\theta^\dagger=1.
\]

\subparagraph*{4. Vacuum injected between the OPAs}

The first column of $U_2$ is

\[
\begin{pmatrix}
c_2\\
e^{-i\phi_2}s_2
\end{pmatrix},
\]

so

\[
\begin{aligned}
F_2(\theta)
&=
\left|
c_2e^{-i\theta}
+s_2e^{i(\theta-\phi_2)}
\right|^2\\
&=c_2^2+s_2^2+2c_2s_2\cos(2\theta-\phi_2)\\
&=\cosh(2r_2)+\sinh(2r_2)\cos(2\theta-\phi_2).
\end{aligned}
\]

\subparagraph*{5. Field propagated through both OPAs}

The product of the two OPA matrices has the form

\[
U_2U_1=
\begin{pmatrix}
A&B\\
B^*&A^*
\end{pmatrix},
\]

where

\[
A=c_2c_1+s_2s_1e^{i(\phi_2-\phi_1)}
\]

and

\[
B=c_2s_1e^{i\phi_1}+s_2c_1e^{i\phi_2}.
\]

Consequently, the first-vacuum contribution is

\[
F_{12}(\theta) = \left| Ae^{-i\theta}+B^*e^{i\theta} \right|^2.
\]

This already gives a compact general answer:

\[
{
\begin{aligned}
\left\langle |I_\theta(\nu)|^2\right\rangle
={}&1-\eta\\
&+\eta(1-\mu)
\left[
\cosh(2r_2)
+\sinh(2r_2)\cos(2\theta-\phi_2)
\right]\\
&+\eta\mu
\left|
Ae^{-i\theta}+B^*e^{i\theta}
\right|^2.
\end{aligned}
}
\]

\subparagraph*{6. Fully expanded general result}

Define

\[
C_j=\cosh(2r_j),\qquad S_j=\sinh(2r_j),
\]

and

\[
\delta=\phi_2-\phi_1,\qquad x=2\theta-\phi_2.
\]

Expanding the last modulus gives

\[
\adjustbox{max width=\linewidth}{$\displaystyle {
\begin{aligned}
\left\langle |I_\theta(\nu)|^2\right\rangle
={}&1-\eta
+\eta(1-\mu)\left(C_2+S_2\cos x\right)\\
&+\eta\mu\big[
C_1C_2+S_1S_2\cos\delta\\
&\qquad\quad+
\left(S_2C_1+S_1C_2\cos\delta\right)\cos x
-S_1\sin\delta\,\sin x
\big].
\end{aligned}
}$}
\]

This expression passes several limiting checks:

\begin{itemize}[leftmargin=*]
\item $\eta=0$ gives vacuum noise, $\langle|I_\theta|^2\rangle=1$.
\item $\mu=0$ leaves only the vacuum amplified by OPA 2.
\item $r_2=0$ gives the ordinary loss-degraded squeezing from OPA 1.
\item $r_1=r_2=0$ gives $\langle|I_\theta|^2\rangle=1$.
\end{itemize}

\paragraph*{Oppositely phased OPAs: $\phi_2-\phi_1=\pi$}

Let

\[
\phi_2=\phi_1+\pi
\]

and define

\[
\chi=2\theta-\phi_1.
\]

The two-OPA matrix simplifies because

\[
A=c_2c_1-s_2s_1=\cosh(r_1-r_2),
\]

and

\[
B=e^{i\phi_1}(c_2s_1-s_2c_1) =e^{i\phi_1}\sinh(r_1-r_2).
\]

Therefore,

\[
F_{12}(\theta) = \cosh\!\left[2(r_1-r_2)\right] + \sinh\!\left[2(r_1-r_2)\right]\cos\chi.
\]

Meanwhile,

\[
F_2(\theta) = \cosh(2r_2)-\sinh(2r_2)\cos\chi.
\]

Hence,

\[
\adjustbox{max width=\linewidth}{$\displaystyle {
\begin{aligned}
\left\langle |I_\theta(\nu)|^2\right\rangle
={}&1-\eta\\
&+\eta(1-\mu)
\left[
\cosh(2r_2)-\sinh(2r_2)\cos\chi
\right]\\
&+\eta\mu
\left[
\cosh\!\left(2(r_1-r_2)\right)
+\sinh\!\left(2(r_1-r_2)\right)\cos\chi
\right].
\end{aligned}
}$}
\]

Equivalently,

\[
\left\langle |I_\theta(\nu)|^2\right\rangle =\mathcal C+\mathcal D\cos\chi,
\]

where

\[
\mathcal C = 1-\eta +\eta\left[ (1-\mu)\cosh(2r_2) +\mu\cosh\!\left(2(r_1-r_2)\right) \right],
\]

and

\[
\mathcal D = \eta\left[ \mu\sinh\!\left(2(r_1-r_2)\right) -(1-\mu)\sinh(2r_2) \right].
\]

Since $-1\leq\cos\chi\leq1$, the minimum squeezed and maximum anti-squeezed powers are

\[
{ V_{\mathrm{sq}}=\mathcal C-|\mathcal D| }
\]

and

\[
{ V_{\mathrm{asq}}=\mathcal C+|\mathcal D|. }
\]

The two orthogonal quadrature values can also be written directly as

\[
V_{\chi=0} = 1-\eta +\eta\left[ \mu e^{2(r_1-r_2)} +(1-\mu)e^{-2r_2} \right],
\]

and

\[
V_{\chi=\pi} = 1-\eta +\eta\left[ \mu e^{-2(r_1-r_2)} +(1-\mu)e^{2r_2} \right].
\]

Thus, without imposing an ordering on $r_1$ and $r_2$,

\[
V_{\mathrm{sq}}=\min\{V_{\chi=0},V_{\chi=\pi}\}, \qquad V_{\mathrm{asq}}=\max\{V_{\chi=0},V_{\chi=\pi}\}.
\]

In the experimentally common readout-amplifier regime $r_2\geq r_1$, $\mathcal D\leq0$. The squeezed quadrature is then $\chi=0$, while the anti-squeezed quadrature is $\chi=\pi$, giving

\[
{ V_{\mathrm{sq}} = 1-\eta +\eta\left[ \mu e^{-2(r_2-r_1)} +(1-\mu)e^{-2r_2} \right] }
\]

and

\[
{ V_{\mathrm{asq}} = 1-\eta +\eta\left[ \mu e^{2(r_2-r_1)} +(1-\mu)e^{2r_2} \right]. }
\]

If $\mu$ were intended as an amplitude rather than a power transmission coefficient, every occurrence of $\mu$ above should be replaced by $\mu^2$.
